\section{Related work}
\label{sec:related}

\subsection{Lateral movement as anomalous-edge detection}

Casting lateral movement as anomalous-edge detection over a temporal graph of
authentication events is an established approach. \textit{Euler}~\cite{euler} composes
a graph neural network with a recurrent encoder over discrete snapshots of a
host-to-host authentication graph and detects red-team activity on LANL~\cite{lanl}.
\textit{Argus}~\cite{argus} reports comparable performance on the same corpus and is
treated as a standard baseline in subsequent literature. \textit{LMDetect}~\cite{lmdetect}
classifies time-aware subgraphs extracted around each authentication event over a
heterogeneous multigraph. A recent \textit{graph foundation model}~\cite{gfm} pretrains
on host-to-host structure and reports results on both LANL and DARPA OpTC with fixed
split conventions (first 41 hours for training on LANL, six days for training and
three for evaluation on OpTC) that have become a standard benchmark protocol.

Reported performance on LANL is governed by the base rate of that corpus: 518 anomalous
edges over 58 days of logs carrying $4.6\times10^{7}$ authentication events. \textit{Euler}
reports an AUC of $0.991$ with an average precision of $0.052$ and a precision of $0.0054$
at its operating point; the graph foundation model reports an AUC of $0.987$ with an average
precision of $0.0088$. Both papers state that AUC is uninformative at this imbalance and
report average precision alongside it. Absolute figures on that corpus are therefore not a
target to match, and do not transfer to a different prevalence or a different decision unit.

One measurement bounds how far the comparison can be taken. \textit{Euler}~\cite{euler}
reports that the ``unknown authentication'' rule, which flags any edge absent from the
training graph, reaches a true positive rate of $72.0\%$ at a false positive rate of $4.4\%$
on LANL. The same rule under the same protocol reaches $63.0\%$ at $42.7\%$ on the stream
used here: the true positive rate is comparable, the false positive rate is an order of
magnitude higher. Benign traffic in this setting also makes non-habitual accesses
(Section~\ref{sec:setup}), so graph memorisation, which accounts for most of the detectable
lateral movement on LANL, is not a shortcut here.

These systems take an authentication log as their substrate: nodes are
hosts (or hosts and users), and edges represent logon events. That formulation fixes what is
detectable. It also explains why they are directly comparable with each other:
LANL carries host-level authentications without client application fingerprints.
LANL is evaluated in Section~\ref{sec:external} as an architectural ablation
(the configuration node removed) rather than as a direct peer comparison.

\subsection{Temporal graph learning}

The architecture builds on the Temporal Graph Network (TGN) of Rossi et al.~\cite{tgn},
which maintains a recurrent per-node memory updated by event messages and computes
embeddings by attending over recent temporal neighbours. The relevant distinction
against the snapshot systems above is temporal granularity. \textit{Euler} and snapshot-based methods discretize the stream into fixed windows and encode snapshot sequences;
a TGN consumes events continuously and updates only the participating entities.
For a PDP that evaluates per request, the continuous formulation directly satisfies
real-time operational constraints, as decisions cannot await window closure. Related continuous
formulations include TGAT~\cite{tgat} and JODIE~\cite{jodie}. TGNs have also
been applied to anomaly detection over dynamic and provenance graphs, though not
over a ZTA request stream with an explicit client-configuration node.

\subsection{Zero Trust and policy-driven authorization}

NIST SP 800-207~\cite{nist800207} separates the policy decision point from the policy
enforcement point and mandates continuous evaluation. In practice the PDP is a policy
engine such as Open Policy Agent~\cite{opa} evaluating declarative rules. A learned
detector in this setting acts as an advisor to that engine, not a replacement.
Consequently, the detector must not unilaterally govern what enters its own baseline,
a requirement enforced via the external commit gate (Section~\ref{sec:problem}).
Surveys of the ZTA literature observe that trust-scoring and PDP proposals
are evaluated predominantly on synthetic simulations, a gap discussed in Section~\ref{sec:external}.

Two recent systems combine Zero Trust with graph neural networks directly.
\textit{TrustGraph}~\cite{trustgraph} builds a heterogeneous graph from windowed
microservice telemetry (logs, metrics, traces, authentication records), scores it
by reconstruction error, and feeds the result through a temporal-decay trust equation
into closed-loop policy enforcement across microservice benchmarks. \textit{ZT-GSAD}~\cite{ztgsad}
applies GraphSAGE under Zero Trust principles to a graph of financial transactions
(accounts, projects, suppliers) for fraud detection. Neither operates on an access-request
stream, includes a node for the client TLS configuration, or evaluates lateral movement
and credential reuse. Additionally, neither implements a commit gate on detector memory updates.
The similarity lies in the high-level framing (Zero Trust combined with graph neural networks),
rather than the underlying substrate, threat model, or streaming protocol.

\subsection{Client fingerprinting}

JA3~\cite{ja3} hashes the ordered parameters of a TLS \texttt{ClientHello} into a
fingerprint of the originating client library and its configuration. Although widely used for
malware-family attribution, JA3 serves here as a \emph{node identity} whose habitual
association with a principal is learned over time. A familiar credential arriving
under an unfamiliar fingerprint thus registers as a structural anomaly rather than a signature match.
This inductive property accommodates previously unseen client tools, provided an adversary
does not replicate the legitimate client TLS parameters.

\subsection{System positioning}

In contrast to prior temporal-graph detectors that consume host-to-host authentication
snapshots without client context, the present system couples a 5-node request chain with
an external commit gate, preserving inductive expressiveness on client configurations while
preventing adversarial drift under streaming execution.
